% Independently written consequences for Illusie I.1.4, printed pp.11-16.
% This file is composed additively with relative-homotopy.tex.

\section{Localization and simplicial derived categories}
\label{section-illusie-I-localization}

\noindent
We use small categories, or a fixed universe enlargement in which all
localizations below exist. A right multiplicative system and its roofs
are as in \href{https://stacks.math.columbia.edu/tag/04VC}{Tag 04VC} and
\href{https://stacks.math.columbia.edu/tag/04VH}{Tag 04VH}.
The following statements spell out the generality and the fixed degree bound
used in Illusie I.1.4; neither requires enough injectives or projectives.

\begin{lemma}[General localization and its derived-functor construction]
\label{lemma-illusie-I-localization-derived}
Let $\mathcal C$ be small and $S$ a collection of its arrows. There is a
functor $Q:\mathcal C\to\mathcal C[S^{-1}]$, identical on objects,
universal for functors sending $S$ to isomorphisms. Natural transformations
between such functors extend uniquely as well.

Suppose now that $S$ is a right multiplicative system and
$F:\mathcal C\to\mathcal B$ is any functor. If for every $X$ the pro-object
$(F(X'))_{s:X'\to X\text{ in }S}$ is isomorphic to a constant object
$L_X$, then these objects define $LF:\mathcal C[S^{-1}]\to\mathcal B$
and a natural map $u:LF\circ Q\to F$ such that
\[
 \operatorname{Nat}(G,LF)\longrightarrow
 \operatorname{Nat}(G\circ Q,F),\qquad a\longmapsto u\circ(aQ),
\]
is bijective for every $G:\mathcal C[S^{-1}]\to\mathcal B$.
The dual assertion holds for a left multiplicative system and essentially
constant ind-objects: $F\to RF\circ Q$ is then universal in the opposite
direction. No additive or triangulated structure is assumed here.
\end{lemma}

\begin{proof}
Adjoin a formal reverse edge $s^{-1}$ for each $s\in S$ to the underlying
graph of $\mathcal C$. Take finite paths and quotient by the congruence
generated by composition and identities in $\mathcal C$, and by
$s^{-1}s=1$ and $ss^{-1}=1$. This constructs the small category
$\mathcal C[S^{-1}]$. The image of every new edge is forced for a functor
inverting $S$, which proves existence and uniqueness of its extension.
Naturality for an invertible arrow implies naturality for its inverse;
hence transformations extend uniquely too.

For the second assertion write $I_X=S/X$, with arbitrary commuting
morphisms over $X$. This category is cofiltered: $1_X$ gives an object,
right Ore squares give common predecessors, and right cancellation
equalizes parallel arrows $a,b$ because $s''a=s''b$ with $s''\in S$.
Let $J_X=(QX\downarrow Q)$. An object of $J_X$
is a pair $(Y,a:QX\to QY)$. The functor
\[
 I_X\longrightarrow J_X,\qquad
 (s:X'\to X)\longmapsto (X',Q(s)^{-1})
\]
is initial. Indeed, objects of its comma category over $(Y,a)$ are precisely
roofs $X\xleftarrow{s}X'\xrightarrow{b}Y$ representing $a$.
There is at least one such roof by right fractions. Any two roofs have a
common refinement giving the same numerator and denominator: first use the
Ore square for the denominators and then the equality criterion for right
fractions (\href{https://stacks.math.columbia.edu/tag/04VJ}{Tag 04VJ}).
Thus this comma category is nonempty and connected, proving initiality.
Equivalently, a compatible cone on $F|_{I_X}$ extends to $F|_{J_X}$ by
choosing a roof for $a$ and composing its cone component with $F(b)$;
common refinements prove independence and compatibility.

An essentially constant pro-object has its representing object as a limit
in $\mathcal B$. Explicitly, full faithfulness of the constant embedding
in the pro-category gives, for every $B$,
\[
 \operatorname{Hom}(B,L_X)
 =\varprojlim_{I_X}\operatorname{Hom}(B,F(X')),
\]
with compatible identifications. Initiality therefore makes $L_X$ a limit
of $F$ over $J_X$. Precomposition by an arrow $QX\to QY$ induces
$J_Y\to J_X$ and hence $L_X\to L_Y$. The uniqueness of limit maps proves
functoriality. Evaluation of the universal cone at $(X,1)$ defines $u_X$.

For $\eta:GQ\to F$, the arrows
$G(QX)\xrightarrow{G(a)}G(QY)\xrightarrow{\eta_Y}F(Y)$
form a cone on $J_X$. Its unique factorization through $L_X$ defines $a_X$;
the same uniqueness proves naturality and $u\circ(aQ)=\eta$.
Conversely, this equation forces every component of the cone and hence
forces $a_X$. This proves the stated bijection. Passing to opposite
categories proves the final assertion.
\end{proof}

\begin{lemma}[Localized Dold--Kan and the fixed nonpositive bound]
\label{lemma-illusie-I-localized-dold-kan}
Let $\mathcal A$ be abelian. Write $C_{\geq0}(\mathcal A)$ for chain
complexes concentrated in nonnegative degrees and $W$ for quasi-isomorphisms.
In $\mathrm{Simp}(\mathcal A)$, let $W_s$ be the maps inducing
quasi-isomorphisms of associated chain complexes. Normalization induces
equivalences
\[
 \mathrm{Simp}(\mathcal A)[W_s^{-1}]
 \simeq C_{\geq0}(\mathcal A)[W^{-1}]
 \simeq D^{\leq0}(\mathcal A),
\]
where $D^{\leq0}(\mathcal A)$ is the full subcategory of $D(\mathcal A)$
with $H^i=0$ for $i>0$, and chains are converted to cochains by $C^i=C_{-i}$.
Both localizations may first be factored through their respective homotopy
categories. The induced equivalences commute with these quotient functors.
Quasi-isomorphisms in $K^{\leq0}(\mathcal A)$ admit both calculi of fractions.
\end{lemma}

\begin{proof}
The natural quasi-isomorphism $N(X)\to s(X)$ of
Lemma \ref{lemma-quasi-isomorphism} shows that $W_s$ is exactly the class
detected by $N$. By Theorem \ref{theorem-dold-kan}, $N$ and its inverse
carry the respective classes to each other. Their units and counits are
isomorphisms, and descend along localization by the preceding lemma.
They remain quasi-inverse, proving the first equivalence.

For clarity the homotopy factorization does not require a calculus of
fractions in the category of complexes itself. For a cochain complex $C$
concentrated in degrees at most zero put
$E^n=C^n\oplus C^{n+1}$, with differential $(a,b)\mapsto(d a+b,-d b)$.
The maps $j:C\to E$, $j(a)=(a,0)$, and
$t:E^n\to E^{n-1}$, $t(a,b)=(0,a)$, satisfy $dt+td=1$.
Thus $E$ is contractible and still concentrated in degrees at most zero.
Set $\operatorname{Cyl}(C)=C\oplus E$, with
$i_0(x)=(x,0)$, $i_1(x)=(x,jx)$, and projection $p$ to $C$.
Since $p$ is a homotopy equivalence and $pi_0=pi_1=1$, any functor
inverting homotopy equivalences identifies $i_0$ and $i_1$.
If $f-g=dh+hd$, the chain map $E\to D$ given by
$(a,b)\mapsto(f-g)(a)+h(b)$, together with $g:C\to D$, extends
$g,f$ along $i_0,i_1$. Hence that functor identifies $f$ and $g$.
Conversely a functor identifying homotopic maps inverts homotopy
equivalences. This proves that passage to the homotopy category is such a
localization. Dold--Kan transports the argument to simplicial objects by
Lemmas \ref{lemma-homotopy-s-N} and \ref{lemma-backwards-homotopy}.
Since homotopy equivalences are quasi-isomorphisms, the asserted
factorizations and commuting comparisons follow.

Let $K^{\leq0}$ be the full homotopy subcategory on complexes concentrated
in degrees at most zero. In $K(\mathcal A)$, quasi-isomorphisms form a
bilateral multiplicative system
(\href{https://stacks.math.columbia.edu/tag/05RT}{Tag 05RT}).
Canonical truncation $\tau_{\leq0}$ is a functor on homotopy categories:
restrict a homotopy in degree zero to $\ker d^0$ and leave its lower
components unchanged. It fixes $K^{\leq0}$ and preserves quasi-isomorphisms.
Apply this functor to the Ore and cancellation diagrams in $K(\mathcal A)$.
All their equalities remain equalities in $K^{\leq0}$ and all required
denominators remain quasi-isomorphisms. This proves the bilateral calculus
claimed above without supposing $K^{\leq0}$ is triangulated.

Every morphism $C\to D$ in $D(\mathcal A)$ between such complexes has a
roof $C\xleftarrow{w}Z\xrightarrow{f}D$. Since $w$ is a
quasi-isomorphism, $H^{>0}(Z)=0$, and $\tau_{\leq0}Z\to Z$ is a
quasi-isomorphism. Truncating the roof represents it inside
$K^{\leq0}[W^{-1}]$, proving fullness. If two roofs there represent the
same arrow in $D(\mathcal A)$, a common refinement in $K(\mathcal A)$
witnesses equality. Its middle complex also has vanishing positive
cohomology. Truncating that refinement witnesses the same equality in
$K^{\leq0}[W^{-1}]$, proving faithfulness. Finally any $Z$ with
$H^{>0}(Z)=0$ is isomorphic to $\tau_{\leq0}Z$, proving that the essential
image is exactly $D^{\leq0}$. This proves the second equivalence.
\end{proof}
